\documentclass{article}
\usepackage{graphicx} % Required for inserting images
\usepackage{amsmath,amssymb,amsxtra,tikz,stackengine,bm,mathrsfs,bbm,tikz-cd,comment,combelow,amsopn}
\usepackage{amsthm}
\usetikzlibrary{arrows.meta}
\usepackage{pgfplots}
\pgfplotsset{compat=1.18}
\usepackage{amscd}
\usepackage{url}

\title{Dehn Twists and Deformation Spaces.}

\begin{document}
\newcommand{\aaa}{\mathfrak{a}}
\newcommand{\cc}{\mathfrak{c}}
\newcommand{\bb}{\mathfrak{b}}
\newcommand{\C}{\mathbb{C}}
\newcommand{\CP}{\mathbb{CP}^1}
\newcommand{\Z}{\mathbb{Z}}
\newcommand{\W}{\mathcal{W}}
\newcommand{\U}{\mathcal{U}}
\newcommand{\V}{\mathcal{V}}
\newcommand{\D}{\mathcal{D}}
\newcommand{\X}{\mathcal{X}}
\newcommand{\Q}{\mathcal{Q}}
\newcommand{\M}{\mathcal{M}}
\newcommand{\Pf}{\mathcal{P}}
\newcommand{\Def}{\mathcal{D}}
\renewcommand{\S}{\mathbb{S}}
\newcommand{\e}{\varepsilon}
\newcommand{\s}{\mathfrak{s}}
\newcommand{\Per}{\operatorname{Per}}
\newcommand{\Mod}{\operatorname{Mod}}
\newcommand{\Mapp}{\operatorname{MCG}}
\newcommand{\Teich}{\operatorname{Teich}}

%\newtheorem*{theorem*}{Theorem}

\newtheorem{theorem}{Theorem}[section]
\newtheorem{lemma}[theorem]{Lemma}
\newtheorem{proposition}[theorem]{Proposition}
\newtheorem{corollary}[theorem]{Corollary}
\newtheorem{conjecture}[theorem]{Conjecture}
\newtheorem{question}[theorem]{Question}

\theoremstyle{definition}
\newtheorem{definition}[theorem]{Definition}
\newtheorem{example}[theorem]{Example}
\newtheorem{construction}[theorem]{Construction}

\theoremstyle{remark}
\newtheorem{remark}[theorem]{Remark}


\numberwithin{equation}{section}
\section{Problem Statement}
Let $f:\S^2\to \S^2$ be a branched cover, $V(f)$ be the set of its critical values,  and $A, B\subset \S^2$ be  finite sets. 
\begin{enumerate}
\item If $f(A)\cup V(f)\subset B$, then every complex structure on $(\S^2, B)$ can be pulled back to a complex structure on $(\S^2, A)$, and this induces the map 
$f^*:\Teich(B)\to \Teich(A)$. 
\item If $A\subset B$, then forgetting the points $B\backslash A$ induces the map $i_*:\Teich(B)\to \Teich(A)$.  
\end{enumerate}
If $f(A)\cup V(f)\cup A\subset B,$ then the deformation space is defined by
$$\D_f:=\{\tau\in \Teich(B): f^*(\tau)=i_*(\tau)\}.$$ 

Consider a quadratic rational map $f$ with critical points $c_1$ and $c_2$. Assume that  $c_1$ belongs to a cycle of length $4$, and the critical value $v=f(c_2)$ does not belong to the cycle. By taking 
$$A=\{c_1, f(c_1), f^{\circ 2}(c_1), f^{\circ 3}(c_1)\}, \quad B=A\cup \{v\},$$ we get a deformation space $\D_f\subset \Teich(B)$. Its image in the moduli space is the set of conjugacy classes of quadratic rational maps with the properties described above. 

Explicitly describe a simple closed curve $\gamma\subset (\S^2, B)$ such that a power of the Dehn twist around this curve preserves the deformation space, and its action on the set of connected components of $\D_f$ has an infinite orbit. 
\end{document}
